\begin{proposition}[the central cusp orbit fibre]\label{push:src:C.3.7}
Let \(H=W/\rho\), and let \(F\cong S^2\) be the fixed double
curve. Then
\[
 \pi_1(H)=\mathbb Z^2,\qquad
 H_i(H;\mathbb Z)=
 \begin{cases}
 \mathbb Z&i=0,3,\\
 \mathbb Z^2&i=1,2,\\
 0&i>3.
 \end{cases}
\]
The class of \(F\) is nonzero and primitive in \(H_2(H;\mathbb Z)\).
In particular \(F\) does not bound a three-chain in the central
orbit fibre. This statement does not assert a deformation retraction
of the entire cusp neighbourhood onto \(H\), or compute the class
of the annular orbit twist.
\end{proposition}
\Uses{the hypotheses, notation and definitions stated here; \cref{push:src:C.3.1}}

\begin{proof}
Number the six toric boundary curves of the normalization
\(\widetilde W=\mathrm{dP}_6\) cyclically as
\(D_1,\ldots,D_6\), with \(D_3,D_6\) the fixed pair, and put
\(p_i=D_i\cap D_{i+1}\), with indices modulo six. Opposite sides
are identified by the maps already checked in Paper~I, Proposition~4.9 and Paper~I, Proposition~4.10 and Paper~I, Theorem~5.10 and \ref{push:src:C.3.1}.
Their endpoint rule is
\[
 p_{i-1}\longmapsto p_{i+3},\qquad
 p_i\longmapsto p_{i+2}\quad(D_i\longrightarrow D_{i+3}).
\]
It follows either from the two triangle classes of the periodic
fan or directly from these six endpoint equations that the vertex
classes are
\(A=\{p_2,p_4,p_6\}\) and \(B=\{p_1,p_3,p_5\}\).
The multiplicative units in the side maps preserve this endpoint
rule. The fixed curve is their resulting sphere, with endpoints
\(A,B\).

Put \(Y=\widetilde W/\rho\). This is an oriented compact
three-manifold whose boundary consists of the two fixed spheres.
Here are the local checks, including the two remaining vertices.
At a free orbit the quotient is an oriented three-dimensional
slice. At a fixed divisor the local model is
\(\mathbb C\times\mathbb C\), with the circle rotating only
the first factor, and its quotient is a half-space. The two toric
vertices not on \(D_3\cup D_6\) are isolated fixed points with
weights \((1,-1)\). Conjugating one coordinate gives the scalar
Hopf action on \(\mathbb C^2\), whose quotient is the cone on
\(\mathbb CP^1=S^2\), namely \(\mathbb R^3\). Thus they are
interior points, not additional boundary components.
All other stabilizers are trivial: the determinants of the circle
cocharacter with the four nonfixed edge rays are \(\pm1\).

The quotient is simply connected. Indeed local slices give path
lifts of orbit-space paths, and lifts on successive subintervals
can be joined by paths in their connected circle orbits. A lifted
loop closes after adding a path in its final orbit; its projection
is the original loop with a constant segment. Since
\(\widetilde W\) is simply connected, the resulting loop is
null-homotopic, and projecting the null-homotopy proves the claim.
Poincar\'e--Lefschetz duality and the boundary sequence now give
\[
 H_1(Y)=H_3(Y)=0,\qquad H_2(Y)=\mathbb Z.
\]
More explicitly, \(H_3(Y,\partial Y)=\mathbb Z\) maps to
\(H_2(\partial Y)=\mathbb Z^2\) by the sum of the two boundary
fundamental classes, and \(H_2(Y,\partial Y)=H^1(Y)=0\).
No classification theorem for simply connected three-manifolds
is needed for these conclusions.

Let \(\widetilde D\subset Y\) be the quotient of the six
boundary curves, and let \(D\subset H\) be its image. The space
\(\widetilde D\) is two spheres joined by two paths, each
subdivided into two edges. Hence
\[
 H_0(\widetilde D)=\mathbb Z,\quad
 H_1(\widetilde D)=\mathbb Z,\quad
 H_2(\widetilde D)=\mathbb Z^2.
\]
The image \(D\) is the sphere \(F\) together with two arcs
\(E_1,E_2\), each joining \(A\) to \(B\). Consequently
\[
 H_0(D)=\mathbb Z,\quad H_1(D)=\mathbb Z^2,\quad
 H_2(D)=\mathbb Z.
\]
Choose a path \(c\subset F\) from \(A\) to \(B\). The
endpoint rule gives the image of the generator of
\(\pi_1(\widetilde D)\) as
\[
 E_1E_2^{-1}cE_1^{-1}E_2c^{-1}.
\]
Writing \(a=E_1c^{-1}\), \(b=E_2c^{-1}\), this is
\(ab^{-1}a^{-1}b\). In particular the induced map on first
homology is zero.

We also fix the second-homology sign. Near a fixed divisor use a
complex normal coordinate \(re^{i\theta}\) and a complex tangent
coordinate. Orient the quotient by putting the positive circle
direction first in the orientation of the surface. At circle
weight \(+1\) the quotient boundary orientation is the complex
curve orientation; at weight \(-1\) it is its negative. The
two normal weights are opposite. Thus, although their boundary
orientations sum to zero in \(H_2(Y)\), the two fixed curves
with their complex orientations represent the same generator
of \(H_2(Y)\). Their maps to \(F\) are holomorphic of degree
one. Therefore the Mayer--Vietoris map
\[
 H_2(\widetilde D)\longrightarrow H_2(Y)\oplus H_2(D)
\]
has matrix
\[
 \begin{pmatrix}1&1\\-1&-1\end{pmatrix}.
\]
It has a primitive rank-one image, rank-one kernel and torsion-free
rank-one cokernel.

The normalization conductor is a topological pushout, and taking
a compact group quotient commutes with that pushout. The inclusion
\(\widetilde D\hookrightarrow Y\) is a cofibration: the finite
toric orbit-stratum decomposition can be triangulated with its
boundary spheres and edge arcs as a subcomplex. Thus this pushout
is a homotopy pushout and its Mayer--Vietoris and van Kampen
sequences apply. Van Kampen gives
\[
 \pi_1(H)=\langle a,b\mid ab^{-1}a^{-1}b=1\rangle
            =\mathbb Z^2.
\]
The preceding homology maps give
\[
 H_3(H)=\mathbb Z,\qquad
 0\longrightarrow\mathbb Z\longrightarrow H_2(H)
      \longrightarrow\mathbb Z\longrightarrow0,
 \qquad H_1(H)=\mathbb Z^2.
\]
The middle sequence splits as a sequence of abelian groups. Its
left generator is the image of \((0,[F])\) in the displayed
second-homology cokernel. It is therefore a primitive nonzero
class in \(H_2(H)\). There are no cells above dimension three.
This proves all assertions, including the nonbounding claim.
The endpoint equivalence classes, incidence ranks, commutator
abelianization, and primitive matrix are checked in
\path{code/d_push_cusp_orbit_cells.py}; the orientation argument
and comparison with the geometric pushout are supplied above.
\end{proof}

